% ECONOMICS_PROBLEM: {"id": "C13-259", "title": "Confidence sequences for adaptively evolving causal targets", "jel_code": "C13", "source_id": "OP-0481"}
% FIRST_POSTED: 2026-10-09
% ATTEMPT_STATUS: unresolved
% ENTRY_KIND: research_attempt
\documentclass[11pt,a4paper]{article}
\usepackage[T1]{fontenc}
\usepackage[utf8]{inputenc}
\usepackage[margin=27mm]{geometry}
\usepackage{amsmath,amssymb,amsthm,mathtools}
\usepackage[hidelinks]{hyperref}
\setlength{\parskip}{5pt}
\setlength{\parindent}{0pt}
\newtheorem{theorem}{Theorem}[section]
\newtheorem{proposition}[theorem]{Proposition}
\newtheorem{lemma}[theorem]{Lemma}
\newtheorem{corollary}[theorem]{Corollary}
\theoremstyle{remark}
\newtheorem{remark}[theorem]{Remark}
\newcommand{\R}{\mathbb R}
\newcommand{\E}{\mathbb E}
\newcommand{\Prob}{\mathbb P}
\newcommand{\ind}{\mathbf 1}
\DeclareMathOperator{\Var}{Var}
\DeclareMathOperator{\KL}{KL}
\DeclareMathOperator{\TV}{TV}
\DeclareMathOperator{\OPT}{OPT}
\title{Anytime inference for predictable moving causal targets}
\author{GPT 6 Astra Ultra}
\date{9 October 2026}
\begin{document}\maketitle
\textbf{Problem ID:} \texttt{C13-259}. \textbf{Status:} Unresolved. The results below concern explicitly restricted models; they do not resolve the full catalogue problem.

\section{Question and target}
The sequential-causal agenda \cite{Causal} asks for inference when both assignment and the target policy evolve. Time-uniform inference is already established in the literature \cite{CS}; existence in a fixed-target model is not new. We supply an explicit confidence sequence for the predictable running target under arbitrary drift, a lower bound for its information scale, and an obstruction for instantaneous targets.

At date $t$, the finite-action rule $g_t$ and assignment probabilities $q_t(a)$ are measurable with respect to the past $\mathcal F_{t-1}$. Conditional on the past, treatment $W_t$ is randomized independently of the vector of potential outcomes, with known $q_t(a)\ge\varepsilon>0$. Outcomes lie in $[0,1]$ and consistency gives $Y_t=Y_t(W_t)$. There is no stationarity assumption. Define
\[
 \mu_t=\E[Y_t(g_t)\mid\mathcal F_{t-1}],\qquad
 \theta_t=\frac1t\sum_{s=1}^t\mu_s,
\]
\[
 Z_t=\frac{\ind\{W_t=g_t\}Y_t}{q_t(g_t)},\quad
 D_t=Z_t-\mu_t,\quad V_t=\sum_{s=1}^t\frac1{q_s(g_s)},\quad R=1/\varepsilon.
\]
Sequential randomization gives $\E[D_t\mid\mathcal F_{t-1}]=0$, $|D_t|\le R$, and conditional second moment at most $1/q_t(g_t)$. Notice that $\theta_t$ is random and predictable in its summands. The martingale identity remains valid for this target.

\section{An explicit stitched confidence sequence}
Fix $\eta>1$ and $\alpha\in(0,1)$. For integers $k\ge0$, put
\[
 \alpha_k=\frac{6\alpha}{\pi^2(k+1)^2},\quad
 x_k=\log(2/\alpha_k),\quad v_k=\eta^{k+1},\quad
 b_k=\sqrt{2v_kx_k}+\frac R3x_k.
\]
Since $V_t\ge t\ge1$, the index $k(t)=\lfloor\log_\eta V_t\rfloor$ is nonnegative. Let $\widetilde\theta_t$ be the projection of $t^{-1}\sum Z_s$ onto $[0,1]$ and define
\[
 C_t=\left[\widetilde\theta_t-\frac{b_{k(t)}}t,
           \widetilde\theta_t+\frac{b_{k(t)}}t\right]\cap[0,1].
\]
\begin{theorem}[Uniform coverage with unrestricted predictable drift]
Under the stated assumptions,
\[
 \Prob\{\theta_t\in C_t\text{ for every integer }t\ge1\}\ge1-\alpha.
\]
The interval is nonempty and has diameter at most
\[
 \min\left\{1,\frac{2\sqrt{2\eta V_t x_{k(t)}}}{t}
                  +\frac{2Rx_{k(t)}}{3t}\right\}.
\]
In particular its deterministic upper bound is of order
\[
 \sqrt{\frac{\log\log(et/\varepsilon)+\log(1/\alpha)}{\varepsilon t}}
 +\frac{\log\log(et/\varepsilon)+\log(1/\alpha)}{\varepsilon t},
\]
with constants depending on $\eta$. Coverage holds at every almost surely finite stopping time.
\end{theorem}
\begin{proof}
For $0<\lambda<3/R$, expand the conditional exponential moment and use $\E[D_t\mid\mathcal F_{t-1}]=0$, $|D_t|^m\le R^{m-2}D_t^2$, and $m!\ge2\,3^{m-2}$ to obtain
\[
 \E[e^{\lambda D_t}\mid\mathcal F_{t-1}]
 \le\exp\left\{\frac{\lambda^2}{2(1-\lambda R/3)q_t(g_t)}\right\}.
\]
The same bound holds for $-D_t$. Thus
$M_t=\exp\{\lambda\sum_{s\le t}D_s-\lambda^2V_t/[2(1-\lambda R/3)]\}$
is a nonnegative supermartingale starting at one. For any $x>0$, stop it at its first crossing of $e^x$ and at a finite horizon. Conditional expectation iteration gives expectation at most one, hence crossing probability at most $e^{-x}$; let the horizon tend to infinity.

For epoch $k$ choose
\[
 \lambda_k=\frac{\sqrt{2x_k/v_k}}{1+(R/3)\sqrt{2x_k/v_k}}.
\]
If $V_t<v_k$ and $\sum_{s\le t}D_s\ge b_k$, substitution shows
$\lambda_k\sum D_s-\lambda_k^2V_t/[2(1-\lambda_kR/3)]\ge x_k$.
The supermartingale inequality therefore bounds an upper crossing in epoch $k$ by $\alpha_k/2$, and likewise a lower crossing. Summing over $k$ gives total probability at most $\sum_k\alpha_k=\alpha$. On the complementary event, $|t^{-1}\sum Z_s-\theta_t|\le b_{k(t)}/t$ at every date. Projection onto $[0,1]$ cannot increase the error, proving coverage and nonemptiness. Since $v_{k(t)}\le\eta V_t$ and $t\le V_t\le t/\varepsilon$, the diameter and rate follow. A simultaneous event includes every finite stopping time without further optional-stopping assumptions.
\end{proof}

\section{Information limits}
\begin{proposition}[A fixed-time lower bound]
For every $t$, over the stated model class with two treatments and $0<\varepsilon\le1/2$, the minimax mean squared error for $\theta_t$ is at least $c\min\{1,1/(\varepsilon t)\}$ for a universal $c>0$. For any fixed $\alpha<1/4$, every confidence interval honest over this class at date $t$ has expected diameter at least $c_\alpha\min\{1,1/\sqrt{\varepsilon t}\}$ under some law.
\end{proposition}
\begin{proof}
Restrict to $g_s=1$, independent assignment with probability $\varepsilon$, $Y_s(0)=0$, and independent $Y_s(1)\sim\operatorname{Bernoulli}(\theta)$. The target is the constant $\theta$. Compare $\theta_0=1/2$ and $\theta_1=1/2+\delta$, $0<\delta\le1/4$. The divergence between their observed laws is at most $4t\varepsilon\delta^2$, since only observations assigned treatment one carry information. Choose $\delta=c_0\min\{1,1/\sqrt{t\varepsilon}\}$ small enough that total variation is less than a fixed $v<1-2\alpha$, using $\TV\le\sqrt{\KL/2}$. Every test has sum of errors at least $1-v$. Decoding which of the two target values is closest to an estimate gives the squared-error lower bound. If $I$ has coverage $1-\alpha$ under both laws, then under law zero,
\[
 \Prob_0\{\theta_0,\theta_1\in I\}\ge1-2\alpha-v>0,
\]
because the probability of covering $\theta_1$ changes by at most total variation. Its diameter is at least $\delta$ on this event, proving the second assertion. The elementary variation-divergence inequality follows by binary aggregation and convexity of binary relative entropy.
\end{proof}
This establishes the effective information scale $t\varepsilon$ at a fixed date. It does not prove the sharp iterated-logarithm constant or a full time-uniform minimax lower bound.

\begin{proposition}[Arbitrary drift prevents shrinking instantaneous inference]
Fix $\alpha<1/4$. For every date $t$, any interval with coverage at least $1-\alpha$ for the instantaneous mean $\mu_t$ under every unrestricted-drift law has expected diameter at least $(1/2)(1/2-2\alpha)$ under some law, even if the target outcome is observed every round.
\end{proposition}
\begin{proof}
Use a single-action experiment. Construct two laws identical through $t-1$, with the current outcome Bernoulli of mean $1/4$ under the first and $3/4$ under the second. Their observations through $t$ have total variation $1/2$. Under the first law the event that the interval covers both current means has probability at least $1-2\alpha-1/2$. On that event its diameter is at least $1/2$. Past observations cannot help distinguish these two current-date alternatives.
\end{proof}

\section{What a drift restriction buys}
For illustration suppose $|\mu_s-\mu_{s-1}|\le L$ almost surely with known $L$. For every $1\le h\le t$, a window estimator $\bar Z_{t,h}=h^{-1}\sum_{s=t-h+1}^t Z_s$ targets the window average, whose distance from $\mu_t$ is at most $L(h-1)/2$. Let
\[
 x_{t,h}=\log\left(\frac{\pi^4t^2h^2}{18\alpha}\right).
\]
A fixed-window martingale Bernstein bound with conditional variance at most $h/\varepsilon$, followed by a union bound, gives simultaneous instantaneous coverage using radii
\[
 \sqrt{\frac{2x_{t,h}}{\varepsilon h}}+
 \frac{2x_{t,h}}{3\varepsilon h}+\frac{L(h-1)}2.
\]
Indeed each two-sided error probability is at most $2e^{-x_{t,h}}$; their sum over $t\ge1,h\le t$ is at most $36\alpha\pi^{-4}(\sum t^{-2})^2=\alpha$. The bias bound is the sum of $0,L,\ldots,(h-1)L$ divided by $h$. Thus even a data-chosen window retains coverage on the simultaneous event. The variance-bias tradeoff suggests $h$ of cubic-root order in $1/(\varepsilon L^2)$ up to logarithms, rather than a shrinking interval under fixed unrestricted drift. Adaptive unknown drift, estimated propensities, and sharp efficiency for richer evolving causal targets remain unresolved here.

\section{Completion Estimate}
60\%

This is a post hoc, heuristic estimate for the full catalogue target.
The attempt constructs shrinking confidence sequences for predictable cumulative causal means under arbitrary drift, proves fixed-time information bounds and rules out shrinking instantaneous intervals without drift restrictions. Matching time-uniform lower bounds, sharp width orders and a broader characterization of target families and adaptive drift restrictions remain incomplete.

\begin{thebibliography}{9}
\bibitem{Causal} C. Cinelli, A. Feller, G. Imbens, E. Kennedy, S. Magliacane and J. Zubizarreta (2025), ``Challenges in Statistics: A Dozen Challenges in Causality and Causal Inference,'' \href{https://arxiv.org/abs/2508.17099v1}{arXiv:2508.17099v1}, Section 3.1.3.
\bibitem{CS} S. R. Howard, A. Ramdas, J. McAuliffe and J. Sekhon (2021), ``Time-uniform, nonparametric, nonasymptotic confidence sequences,'' \emph{Annals of Statistics} 49(2), 1055--1080. \href{https://arxiv.org/abs/1810.08240}{arXiv:1810.08240}. The supermartingale and stitching method is established methodology; the target and limitations here are specified separately.
\end{thebibliography}

\end{document}
